\documentclass[10pt,a4paper]{amsart}
\usepackage[margin=19mm]{geometry}
\usepackage{amsmath,amssymb,mathtools}
\usepackage[scr=rsfs,cal=euler]{mathalfa}
\usepackage{xfrac}
\usepackage[unicode,hidelinks,bookmarks=false]{hyperref}
\newcommand{\ie}{i.e.\ }
\newcommand{\cf}{cf.\ }
\newcommand{\resp}{respectively\ }
\newcommand{\Subsection}{\subsection}
\providecommand{\nobreakdash}{\nobreak\hbox{-}}
\setlength{\emergencystretch}{2em}
\multlinegap0pt
\usepackage{amssymb}
\newtheorem{theorem}{Theorem}
\newtheorem{corollary}{Corollary}
\title{Singular inner eigenfunctions of composition operators}
\author{V. A. Anjali, P. Muthukumar, P. Shankar}
\date{Source: arXiv:2608.29722v1 (CC BY 4.0)}
\begin{document}
\maketitle

\noindent\emph{Source and licence.} Singular inner eigenfunctions of composition operators. Version arXiv:2608.29722v1 (CC BY 4.0), distributed by arXiv under the Creative Commons Attribution 4.0 International licence, http://creativecommons.org/licenses/by/4.0/, as recorded in the arXiv licence metadata for this exact version and shown on the abstract page https://arxiv.org/abs/2608.29722v1. Full licence notice: \texttt{licenses/arxiv-2608.29722-CC-BY-4.0.txt}.\par\smallskip

\noindent\textit{Editorial scope.} Counted unit: the article's corollary corollary (source0.tex lines 628--637). The article's complete original proof of it is delivered with it (source0.tex lines 638--647, one \texttt{proof} environment of 10 lines). No mathematics is written, repaired or re-derived: the statement and the proof are the article's own lines, in the article's own notation, and no retained line is substituted or re-flowed.  Retained uncounted as the source-local context the proof needs: lines 616--623.  Nothing was omitted from any retained span, so the package carries no omission record.  No retained line contains a citation, so the wrapper carries no bibliography and no bibliography entry.  No retained line contains a figure environment, an image-inclusion command or drawing code, so no image is delivered and the package declares no asset dependency.  Editorial formatting only, in addition: an AMS \texttt{amsart} preamble that restates the article's own declarations and macros used by the retained lines, namely the environments \texttt{theorem}, \texttt{corollary} on separate counters, together with the standard packages those lines need.  The retained environments print in this document as Theorem 1, Corollary 1 in order of appearance, whereas the article prints the same items with the numbering of its own full document.  The complete native source of the article as distributed in the arXiv e-print for this exact version is bundled unchanged as \texttt{sources/208883/source0.tex}.
\begin{theorem}\label{pure}
	Let $\phi$ be a non-elliptic automorphism of $\mathbb{D}$ and $\mu$ be a discrete singular measure on $\mathbb{T}$. Then $C_{\phi}(zS_{\mu}H^2)^{\perp}\subseteq(zS_{\mu}H^2)^{\perp}$
	if and only if 
	\begin{enumerate}
		\item $\phi^{-1}(atom(\mu))\subseteq{atom(\mu)}$.
		\item $\dfrac{\mu\{\phi^{-1}(w)\}}{\mu\{w\}}\geq |({\phi^{-1}})^{'}(w)|$ for all $w \in atom{(\mu)}$.
	\end{enumerate}
\end{theorem}

\begin{corollary}
	Let $S(z)=\exp(\alpha\frac{z+p}{z-p})$ be an atomic singular inner function, where  $\alpha>0$ and $p \in \mathbb{T}$.
	\begin{enumerate}
		\item	Let $\phi$  be a parabolic automorphism. Then $C_\phi(zSH^2)^\perp\subseteq (zSH^2)^\perp$ if and only if
		$p$ is the fixed point of $\phi$.
		\item  Let $\phi$  be a hyperbolic automorphism. Then $C_\phi(zSH^2)^\perp\subseteq (zSH^2)^\perp$ if and only if
		$p$ the fixed point of $\phi$ and $p$ is not the Denjoy-Wolff point.
\end{enumerate}
	
\end{corollary} 

\begin{proof}
	
Proof of (1):	Let $\phi$ be a parabolic automorphism.
	If $C_\phi(zSH^2)^\perp\subseteq (zSH^2)^\perp$ , then by the first condition of Theorem \ref{pure}, we have $p \in \{\phi(p)\}$ and thus $p$ is the fixed point of $\phi$. The converse part directly follows from Theorem \ref{pure}.
	% We have $atom({\mu})=\{p\}$.
	 
Proof of (2):	Let $\phi$ be a hyperbolic automorphism. Assume that $C_\phi(zSH^2)^\perp\subseteq (zSH^2)^\perp$. Then condition (1) in Theorem \ref{pure} implies $\phi(p)=p$. We can see that, for any $w \in \mathbb{T}$, $(\phi^{-1})^{'}(w)=\frac{1}{\phi^{'}(\phi^{-1}(w))}$. By condition (2) in Theorem \ref{pure}, we have 
$|\phi^{'}(p)|\geq 1$, which is only possible if $p$ is not the  Denjoy–Wolff point. The converse part directly follows from Theorem \ref{pure}.

\end{proof}

\end{document}
